% --- Archon physics formalization source begin ---
% archon:physics
% archon:covers PhyXMiniProblems/problem_phyx_mini_0626.lean
% archon:source-report reports/phyx_mini/problem_phyx_mini_0626.source.json

\chapter{Physics problem phyx\_mini\_0626}
\label{ch:PhyXMiniProblems_problem_phyx_mini_0626}

\paragraph{Problem source.}
\#\# Physical scenario

A strip of silicon \textbackslash{}(1.6\textbackslash{},\textbackslash{}mathrm\{cm\}\textbackslash{}) wide and \textbackslash{}(1.0\textbackslash{},\textbackslash{}mathrm\{mm\}\textbackslash{}) thick is immersed in a magnetic field of strength \textbackslash{}(1.5\textbackslash{},\textbackslash{}mathrm\{T\}\textbackslash{}) perpendicular to the strip (figure). When a current of \textbackslash{}(0.28\textbackslash{},\textbackslash{}mathrm\{mA\}\textbackslash{}) is run through the strip, there is a resulting Hall effect voltage of \textbackslash{}(18\textbackslash{},\textbackslash{}mathrm\{mV\}\textbackslash{}) across the strip. The density of silicon is \textbackslash{}(2330\textbackslash{},\textbackslash{}mathrm\{kg/m\textasciicircum{}3\}\textbackslash{}).

\#\# Question

How many electrons per silicon atom are in the conduction band?

\#\# Answer choices

- A: A: \$3.9\textbackslash{}times10\textasciicircum{}\{-9\}\textbackslash{}

- B: B: \$2.9\textbackslash{}times10\textasciicircum{}\{-8\}\textbackslash{}

- C: C: \$2.0\textbackslash{}times10\textasciicircum{}\{-9\}\textbackslash{}

- D: D: \$2.9\textbackslash{}times10\textasciicircum{}\{-9\}\textbackslash{}

\#\# Figure caption (auxiliary; use the image as primary evidence)

The image depicts a rectangular, flat surface with two labeled vectors. 

1. **Surface**: The object appears to be a flat, rectangular slab oriented at an angle for a three-dimensional perspective.

2. **Vectors**:
   - **Current (I)**: A black arrow labeled "I" points along the surface's lower edge, indicating the direction of the current flow.
   - **Magnetic Field (B)**: A blue arrow labeled with "B" points perpendicular to the surface, representing a magnetic field.

3. **Labels and Direction**:
   - The vector labeled "I" has a leftward direction, staying parallel to the surface plane.
   - The vector labeled "B" points upward and away from the surface, indicating that the magnetic field is perpendicular to the plane of the surface. 

4. **Relationships**:
   - The presence of both vectors suggests a scenario involving electromagnetic properties, such as the Hall effect, where a current and magnetic field interact perpendicularly.

\#\# Dataset metadata

Quantum Phenomena; Multi-Formula Reasoning, Physical Model Grounding Reasoning

\paragraph{Recorded answer/context.}
D: D: \$2.9\textbackslash{}times10\textasciicircum{}\{-9\}\textbackslash{}

\paragraph{Figure/image path.}
/root/proposal\_for\_physic/science-mango/phyx\_mini\_run/phyx\_data/test\_image/626.png

\paragraph{Formalization target.}
create a compiling Lean file with sorry bodies at `PhyXMiniProblems/problem\_phyx\_mini\_0626.lean`. The Lean declarations must preserve the physical quantities, dimensions or dimensional roles, figure labels, governing-law hypotheses, and final relation expressed by this problem.
Use Mathlib/Physlib names found through LeanExplore where available. If a physics API is missing, introduce faithful local abstractions rather than scalar placeholder aliases.

\begin{theorem}[Physics formalization target]
\label{thm:physics:phyx_mini_0626:target}
\lean{PhyXMiniProblems.ProblemPhyXMini0626.conductionElectronsPerSiliconAtom_is_answer_D}
\uses{def:physics:phyx-mini-0626:phyxminiproblems-problemphyxmini0626-siliconhallsetup, def:physics:phyx-mini-0626:phyxminiproblems-problemphyxmini0626-matchessiliconhallscenario, def:physics:phyx-mini-0626:phyxminiproblems-problemphyxmini0626-matchessuppliedhallstripfigure, def:physics:phyx-mini-0626:phyxminiproblems-problemphyxmini0626-matchesproblemmeasurements, def:physics:phyx-mini-0626:phyxminiproblems-problemphyxmini0626-matchessiliconatomicreferencedata, def:physics:phyx-mini-0626:phyxminiproblems-problemphyxmini0626-satisfieshalleffectlaw, def:physics:phyx-mini-0626:phyxminiproblems-problemphyxmini0626-satisfiessiliconatomicdensitylaw, def:physics:phyx-mini-0626:phyxminiproblems-problemphyxmini0626-satisfiescarrierperatombalance, def:physics:phyx-mini-0626:phyxminiproblems-problemphyxmini0626-answerchoice, def:physics:phyx-mini-0626:phyxminiproblems-problemphyxmini0626-isclosestanswer, def:physics:phyx-mini-0626:phyxminiproblems-problemphyxmini0626-roundstodisplayedvalue}
The assigned autoformalize agent should translate this physics problem into Lean declarations in the covered file, with theorem and lemma proof bodies written as `by sorry`.
\end{theorem}
\begin{proof}
This is an autoformalization task, not a proof task. Produce statements that can later be proved by the physics prover without weakening the physical model.
\end{proof}

% --- Archon declaration topology begin ---
\section{Declaration topology}
\begin{definition}[electric Current Dimension]
  \label{def:physics:phyx-mini-0626:phyxminiproblems-problemphyxmini0626-electriccurrentdimension}
  \lean{PhyXMiniProblems.ProblemPhyXMini0626.electricCurrentDimension}
  Electric current has dimension charge per time
\end{definition}
\begin{proof}
  This is the declared model definition of electric Current Dimension.
\end{proof}
\begin{definition}[energy Dimension]
  \label{def:physics:phyx-mini-0626:phyxminiproblems-problemphyxmini0626-energydimension}
  \lean{PhyXMiniProblems.ProblemPhyXMini0626.energyDimension}
  Energy in the MLTQ dimension system
\end{definition}
\begin{proof}
  This is the declared model definition of energy Dimension.
\end{proof}
\begin{definition}[electric Potential Dimension]
  \label{def:physics:phyx-mini-0626:phyxminiproblems-problemphyxmini0626-electricpotentialdimension}
  \lean{PhyXMiniProblems.ProblemPhyXMini0626.electricPotentialDimension}
  \uses{def:physics:phyx-mini-0626:phyxminiproblems-problemphyxmini0626-energydimension}
  Electric potential has dimension energy per charge
\end{definition}
\begin{proof}
  This is the declared model definition of electric Potential Dimension.
\end{proof}
\begin{definition}[magnetic Flux Density Dimension]
  \label{def:physics:phyx-mini-0626:phyxminiproblems-problemphyxmini0626-magneticfluxdensitydimension}
  \lean{PhyXMiniProblems.ProblemPhyXMini0626.magneticFluxDensityDimension}
  Magnetic flux density (tesla) has dimension mass per charge per time
\end{definition}
\begin{proof}
  This is the declared model definition of magnetic Flux Density Dimension.
\end{proof}
\begin{definition}[mass Density Dimension]
  \label{def:physics:phyx-mini-0626:phyxminiproblems-problemphyxmini0626-massdensitydimension}
  \lean{PhyXMiniProblems.ProblemPhyXMini0626.massDensityDimension}
  Bulk mass density has dimension mass per volume
\end{definition}
\begin{proof}
  This is the declared model definition of mass Density Dimension.
\end{proof}
\begin{definition}[number Density Dimension]
  \label{def:physics:phyx-mini-0626:phyxminiproblems-problemphyxmini0626-numberdensitydimension}
  \lean{PhyXMiniProblems.ProblemPhyXMini0626.numberDensityDimension}
  Particle number density has dimension inverse volume
\end{definition}
\begin{proof}
  This is the declared model definition of number Density Dimension.
\end{proof}
\begin{definition}[Length Quantity]
  \label{def:physics:phyx-mini-0626:phyxminiproblems-problemphyxmini0626-lengthquantity}
  \lean{PhyXMiniProblems.ProblemPhyXMini0626.LengthQuantity}
  A nonnegative physical length
\end{definition}
\begin{proof}
  This is the declared model definition of Length Quantity.
\end{proof}
\begin{definition}[Electric Current Quantity]
  \label{def:physics:phyx-mini-0626:phyxminiproblems-problemphyxmini0626-electriccurrentquantity}
  \lean{PhyXMiniProblems.ProblemPhyXMini0626.ElectricCurrentQuantity}
  \uses{def:physics:phyx-mini-0626:phyxminiproblems-problemphyxmini0626-electriccurrentdimension}
  A nonnegative physical electric-current magnitude
\end{definition}
\begin{proof}
  This is the declared model definition of Electric Current Quantity.
\end{proof}
\begin{definition}[Electric Potential Quantity]
  \label{def:physics:phyx-mini-0626:phyxminiproblems-problemphyxmini0626-electricpotentialquantity}
  \lean{PhyXMiniProblems.ProblemPhyXMini0626.ElectricPotentialQuantity}
  \uses{def:physics:phyx-mini-0626:phyxminiproblems-problemphyxmini0626-electricpotentialdimension}
  A nonnegative physical Hall-voltage magnitude
\end{definition}
\begin{proof}
  This is the declared model definition of Electric Potential Quantity.
\end{proof}
\begin{definition}[Magnetic Flux Density Quantity]
  \label{def:physics:phyx-mini-0626:phyxminiproblems-problemphyxmini0626-magneticfluxdensityquantity}
  \lean{PhyXMiniProblems.ProblemPhyXMini0626.MagneticFluxDensityQuantity}
  \uses{def:physics:phyx-mini-0626:phyxminiproblems-problemphyxmini0626-magneticfluxdensitydimension}
  A nonnegative physical magnetic-flux-density magnitude
\end{definition}
\begin{proof}
  This is the declared model definition of Magnetic Flux Density Quantity.
\end{proof}
\begin{definition}[Charge Magnitude Quantity]
  \label{def:physics:phyx-mini-0626:phyxminiproblems-problemphyxmini0626-chargemagnitudequantity}
  \lean{PhyXMiniProblems.ProblemPhyXMini0626.ChargeMagnitudeQuantity}
  A nonnegative physical charge magnitude
\end{definition}
\begin{proof}
  This is the declared model definition of Charge Magnitude Quantity.
\end{proof}
\begin{definition}[Mass Quantity]
  \label{def:physics:phyx-mini-0626:phyxminiproblems-problemphyxmini0626-massquantity}
  \lean{PhyXMiniProblems.ProblemPhyXMini0626.MassQuantity}
  A nonnegative physical mass
\end{definition}
\begin{proof}
  This is the declared model definition of Mass Quantity.
\end{proof}
\begin{definition}[Mass Density Quantity]
  \label{def:physics:phyx-mini-0626:phyxminiproblems-problemphyxmini0626-massdensityquantity}
  \lean{PhyXMiniProblems.ProblemPhyXMini0626.MassDensityQuantity}
  \uses{def:physics:phyx-mini-0626:phyxminiproblems-problemphyxmini0626-massdensitydimension}
  A nonnegative physical bulk mass density
\end{definition}
\begin{proof}
  This is the declared model definition of Mass Density Quantity.
\end{proof}
\begin{definition}[Number Density Quantity]
  \label{def:physics:phyx-mini-0626:phyxminiproblems-problemphyxmini0626-numberdensityquantity}
  \lean{PhyXMiniProblems.ProblemPhyXMini0626.NumberDensityQuantity}
  \uses{def:physics:phyx-mini-0626:phyxminiproblems-problemphyxmini0626-numberdensitydimension}
  A nonnegative physical particle number density
\end{definition}
\begin{proof}
  This is the declared model definition of Number Density Quantity.
\end{proof}
\begin{definition}[Spatial Direction]
  \label{def:physics:phyx-mini-0626:phyxminiproblems-problemphyxmini0626-spatialdirection}
  \lean{PhyXMiniProblems.ProblemPhyXMini0626.SpatialDirection}
  A direction vector in physical three-space
\end{definition}
\begin{proof}
  This is the declared model definition of Spatial Direction.
\end{proof}
\begin{definition}[nonnegative SIReadout]
  \label{def:physics:phyx-mini-0626:phyxminiproblems-problemphyxmini0626-nonnegativesireadout}
  \lean{PhyXMiniProblems.ProblemPhyXMini0626.nonnegativeSIReadout}
  Coherent-SI readout of a nonnegative dimensionful quantity
\end{definition}
\begin{proof}
  This is the declared model definition of nonnegative SIReadout.
\end{proof}
\begin{definition}[length Readout]
  \label{def:physics:phyx-mini-0626:phyxminiproblems-problemphyxmini0626-lengthreadout}
  \lean{PhyXMiniProblems.ProblemPhyXMini0626.lengthReadout}
  \uses{def:physics:phyx-mini-0626:phyxminiproblems-problemphyxmini0626-lengthquantity}
  Read a physical length in a selected Physlib length unit
\end{definition}
\begin{proof}
  This is the declared model definition of length Readout.
\end{proof}
\begin{definition}[electric Current In Amperes]
  \label{def:physics:phyx-mini-0626:phyxminiproblems-problemphyxmini0626-electriccurrentinamperes}
  \lean{PhyXMiniProblems.ProblemPhyXMini0626.electricCurrentInAmperes}
  \uses{def:physics:phyx-mini-0626:phyxminiproblems-problemphyxmini0626-electriccurrentquantity, def:physics:phyx-mini-0626:phyxminiproblems-problemphyxmini0626-nonnegativesireadout}
  Electric-current readout in amperes (coulomb / second)
\end{definition}
\begin{proof}
  This is the declared model definition of electric Current In Amperes.
\end{proof}
\begin{definition}[electric Current In Milliamperes]
  \label{def:physics:phyx-mini-0626:phyxminiproblems-problemphyxmini0626-electriccurrentinmilliamperes}
  \lean{PhyXMiniProblems.ProblemPhyXMini0626.electricCurrentInMilliamperes}
  \uses{def:physics:phyx-mini-0626:phyxminiproblems-problemphyxmini0626-electriccurrentquantity, def:physics:phyx-mini-0626:phyxminiproblems-problemphyxmini0626-electriccurrentinamperes}
  Electric-current readout in milliamperes
\end{definition}
\begin{proof}
  This is the declared model definition of electric Current In Milliamperes.
\end{proof}
\begin{definition}[electric Potential In Volts]
  \label{def:physics:phyx-mini-0626:phyxminiproblems-problemphyxmini0626-electricpotentialinvolts}
  \lean{PhyXMiniProblems.ProblemPhyXMini0626.electricPotentialInVolts}
  \uses{def:physics:phyx-mini-0626:phyxminiproblems-problemphyxmini0626-electricpotentialquantity, def:physics:phyx-mini-0626:phyxminiproblems-problemphyxmini0626-nonnegativesireadout}
  Hall-voltage-magnitude readout in volts
\end{definition}
\begin{proof}
  This is the declared model definition of electric Potential In Volts.
\end{proof}
\begin{definition}[electric Potential In Millivolts]
  \label{def:physics:phyx-mini-0626:phyxminiproblems-problemphyxmini0626-electricpotentialinmillivolts}
  \lean{PhyXMiniProblems.ProblemPhyXMini0626.electricPotentialInMillivolts}
  \uses{def:physics:phyx-mini-0626:phyxminiproblems-problemphyxmini0626-electricpotentialquantity, def:physics:phyx-mini-0626:phyxminiproblems-problemphyxmini0626-electricpotentialinvolts}
  Hall-voltage-magnitude readout in millivolts
\end{definition}
\begin{proof}
  This is the declared model definition of electric Potential In Millivolts.
\end{proof}
\begin{definition}[magnetic Flux Density In Teslas]
  \label{def:physics:phyx-mini-0626:phyxminiproblems-problemphyxmini0626-magneticfluxdensityinteslas}
  \lean{PhyXMiniProblems.ProblemPhyXMini0626.magneticFluxDensityInTeslas}
  \uses{def:physics:phyx-mini-0626:phyxminiproblems-problemphyxmini0626-magneticfluxdensityquantity, def:physics:phyx-mini-0626:phyxminiproblems-problemphyxmini0626-nonnegativesireadout}
  Magnetic-flux-density readout in teslas
\end{definition}
\begin{proof}
  This is the declared model definition of magnetic Flux Density In Teslas.
\end{proof}
\begin{definition}[charge Magnitude In Coulombs]
  \label{def:physics:phyx-mini-0626:phyxminiproblems-problemphyxmini0626-chargemagnitudeincoulombs}
  \lean{PhyXMiniProblems.ProblemPhyXMini0626.chargeMagnitudeInCoulombs}
  \uses{def:physics:phyx-mini-0626:phyxminiproblems-problemphyxmini0626-chargemagnitudequantity, def:physics:phyx-mini-0626:phyxminiproblems-problemphyxmini0626-nonnegativesireadout}
  Charge-magnitude readout in coulombs
\end{definition}
\begin{proof}
  This is the declared model definition of charge Magnitude In Coulombs.
\end{proof}
\begin{definition}[charge Magnitude In Elementary Charges]
  \label{def:physics:phyx-mini-0626:phyxminiproblems-problemphyxmini0626-chargemagnitudeinelementarycharges}
  \lean{PhyXMiniProblems.ProblemPhyXMini0626.chargeMagnitudeInElementaryCharges}
  \uses{def:physics:phyx-mini-0626:phyxminiproblems-problemphyxmini0626-chargemagnitudequantity}
  Charge-magnitude readout in units of Physlib's elementary charge
\end{definition}
\begin{proof}
  This is the declared model definition of charge Magnitude In Elementary Charges.
\end{proof}
\begin{definition}[mass In Kilograms]
  \label{def:physics:phyx-mini-0626:phyxminiproblems-problemphyxmini0626-massinkilograms}
  \lean{PhyXMiniProblems.ProblemPhyXMini0626.massInKilograms}
  \uses{def:physics:phyx-mini-0626:phyxminiproblems-problemphyxmini0626-massquantity, def:physics:phyx-mini-0626:phyxminiproblems-problemphyxmini0626-nonnegativesireadout}
  Mass readout in kilograms
\end{definition}
\begin{proof}
  This is the declared model definition of mass In Kilograms.
\end{proof}
\begin{definition}[mass Density In Kilograms Per Cubic Meter]
  \label{def:physics:phyx-mini-0626:phyxminiproblems-problemphyxmini0626-massdensityinkilogramspercubicmeter}
  \lean{PhyXMiniProblems.ProblemPhyXMini0626.massDensityInKilogramsPerCubicMeter}
  \uses{def:physics:phyx-mini-0626:phyxminiproblems-problemphyxmini0626-massdensityquantity, def:physics:phyx-mini-0626:phyxminiproblems-problemphyxmini0626-nonnegativesireadout}
  Bulk-mass-density readout in kilograms per cubic metre
\end{definition}
\begin{proof}
  This is the declared model definition of mass Density In Kilograms Per Cubic Meter.
\end{proof}
\begin{definition}[number Density Per Cubic Meter]
  \label{def:physics:phyx-mini-0626:phyxminiproblems-problemphyxmini0626-numberdensitypercubicmeter}
  \lean{PhyXMiniProblems.ProblemPhyXMini0626.numberDensityPerCubicMeter}
  \uses{def:physics:phyx-mini-0626:phyxminiproblems-problemphyxmini0626-numberdensityquantity, def:physics:phyx-mini-0626:phyxminiproblems-problemphyxmini0626-nonnegativesireadout}
  Particle-number-density readout in particles per cubic metre
\end{definition}
\begin{proof}
  This is the declared model definition of number Density Per Cubic Meter.
\end{proof}
\begin{definition}[Figure Vector]
  \label{def:physics:phyx-mini-0626:phyxminiproblems-problemphyxmini0626-figurevector}
  \lean{PhyXMiniProblems.ProblemPhyXMini0626.FigureVector}
  The two vector labels visible in the supplied image
\end{definition}
\begin{proof}
  This is the declared model definition of Figure Vector.
\end{proof}
\begin{definition}[Figure Arrow Color]
  \label{def:physics:phyx-mini-0626:phyxminiproblems-problemphyxmini0626-figurearrowcolor}
  \lean{PhyXMiniProblems.ProblemPhyXMini0626.FigureArrowColor}
  The visible color of a vector arrow in the supplied image
\end{definition}
\begin{proof}
  This is the declared model definition of Figure Arrow Color.
\end{proof}
\begin{definition}[Hall Strip Figure]
  \label{def:physics:phyx-mini-0626:phyxminiproblems-problemphyxmini0626-hallstripfigure}
  \lean{PhyXMiniProblems.ProblemPhyXMini0626.HallStripFigure}
  \uses{def:physics:phyx-mini-0626:phyxminiproblems-problemphyxmini0626-figurevector, def:physics:phyx-mini-0626:phyxminiproblems-problemphyxmini0626-figurearrowcolor}
  Primary-image information from 626.png
\end{definition}
\begin{proof}
  This is the declared model definition of Hall Strip Figure.
\end{proof}
\begin{definition}[Silicon Hall Setup]
  \label{def:physics:phyx-mini-0626:phyxminiproblems-problemphyxmini0626-siliconhallsetup}
  \lean{PhyXMiniProblems.ProblemPhyXMini0626.SiliconHallSetup}
  \uses{def:physics:phyx-mini-0626:phyxminiproblems-problemphyxmini0626-lengthquantity, def:physics:phyx-mini-0626:phyxminiproblems-problemphyxmini0626-electriccurrentquantity, def:physics:phyx-mini-0626:phyxminiproblems-problemphyxmini0626-electricpotentialquantity, def:physics:phyx-mini-0626:phyxminiproblems-problemphyxmini0626-magneticfluxdensityquantity, def:physics:phyx-mini-0626:phyxminiproblems-problemphyxmini0626-chargemagnitudequantity, def:physics:phyx-mini-0626:phyxminiproblems-problemphyxmini0626-massquantity, def:physics:phyx-mini-0626:phyxminiproblems-problemphyxmini0626-massdensityquantity, def:physics:phyx-mini-0626:phyxminiproblems-problemphyxmini0626-numberdensityquantity, def:physics:phyx-mini-0626:phyxminiproblems-problemphyxmini0626-spatialdirection, def:physics:phyx-mini-0626:phyxminiproblems-problemphyxmini0626-hallstripfigure}
  The physical silicon-strip experiment. The requested carrier ratio is an independent dimensionless observable; it is not defined from an answer choice or from the numerical result to be shown
\end{definition}
\begin{proof}
  This is the declared model definition of Silicon Hall Setup.
\end{proof}
\begin{definition}[Matches Silicon Hall Scenario]
  \label{def:physics:phyx-mini-0626:phyxminiproblems-problemphyxmini0626-matchessiliconhallscenario}
  \lean{PhyXMiniProblems.ProblemPhyXMini0626.MatchesSiliconHallScenario}
  \uses{def:physics:phyx-mini-0626:phyxminiproblems-problemphyxmini0626-lengthreadout, def:physics:phyx-mini-0626:phyxminiproblems-problemphyxmini0626-electriccurrentinamperes, def:physics:phyx-mini-0626:phyxminiproblems-problemphyxmini0626-electricpotentialinvolts, def:physics:phyx-mini-0626:phyxminiproblems-problemphyxmini0626-magneticfluxdensityinteslas, def:physics:phyx-mini-0626:phyxminiproblems-problemphyxmini0626-numberdensitypercubicmeter, def:physics:phyx-mini-0626:phyxminiproblems-problemphyxmini0626-siliconhallsetup}
  Qualitative geometry and positivity conditions of the Hall experiment
\end{definition}
\begin{proof}
  This is the declared model definition of Matches Silicon Hall Scenario.
\end{proof}
\begin{definition}[Matches Supplied Hall Strip Figure]
  \label{def:physics:phyx-mini-0626:phyxminiproblems-problemphyxmini0626-matchessuppliedhallstripfigure}
  \lean{PhyXMiniProblems.ProblemPhyXMini0626.MatchesSuppliedHallStripFigure}
  \uses{def:physics:phyx-mini-0626:phyxminiproblems-problemphyxmini0626-hallstripfigure}
  Labels, colors, and arrow geometry read directly from the supplied image
\end{definition}
\begin{proof}
  This is the declared model definition of Matches Supplied Hall Strip Figure.
\end{proof}
\begin{definition}[Matches Problem Measurements]
  \label{def:physics:phyx-mini-0626:phyxminiproblems-problemphyxmini0626-matchesproblemmeasurements}
  \lean{PhyXMiniProblems.ProblemPhyXMini0626.MatchesProblemMeasurements}
  \uses{def:physics:phyx-mini-0626:phyxminiproblems-problemphyxmini0626-lengthreadout, def:physics:phyx-mini-0626:phyxminiproblems-problemphyxmini0626-electriccurrentinmilliamperes, def:physics:phyx-mini-0626:phyxminiproblems-problemphyxmini0626-electricpotentialinmillivolts, def:physics:phyx-mini-0626:phyxminiproblems-problemphyxmini0626-magneticfluxdensityinteslas, def:physics:phyx-mini-0626:phyxminiproblems-problemphyxmini0626-massdensityinkilogramspercubicmeter, def:physics:phyx-mini-0626:phyxminiproblems-problemphyxmini0626-siliconhallsetup}
  Measured dimensions and electromagnetic readouts stated in the problem
\end{definition}
\begin{proof}
  This is the declared model definition of Matches Problem Measurements.
\end{proof}
\begin{definition}[Matches Silicon Atomic Reference Data]
  \label{def:physics:phyx-mini-0626:phyxminiproblems-problemphyxmini0626-matchessiliconatomicreferencedata}
  \lean{PhyXMiniProblems.ProblemPhyXMini0626.MatchesSiliconAtomicReferenceData}
  \uses{def:physics:phyx-mini-0626:phyxminiproblems-problemphyxmini0626-chargemagnitudeinelementarycharges, def:physics:phyx-mini-0626:phyxminiproblems-problemphyxmini0626-siliconhallsetup}
  Reference atomic data needed to convert the stated silicon mass density to an atom number density. These are calibrated data readouts, not the requested conduction-electron ratio
\end{definition}
\begin{proof}
  This is the declared model definition of Matches Silicon Atomic Reference Data.
\end{proof}
\begin{definition}[Satisfies Hall Effect Law]
  \label{def:physics:phyx-mini-0626:phyxminiproblems-problemphyxmini0626-satisfieshalleffectlaw}
  \lean{PhyXMiniProblems.ProblemPhyXMini0626.SatisfiesHallEffectLaw}
  \uses{def:physics:phyx-mini-0626:phyxminiproblems-problemphyxmini0626-lengthreadout, def:physics:phyx-mini-0626:phyxminiproblems-problemphyxmini0626-electriccurrentinamperes, def:physics:phyx-mini-0626:phyxminiproblems-problemphyxmini0626-electricpotentialinvolts, def:physics:phyx-mini-0626:phyxminiproblems-problemphyxmini0626-magneticfluxdensityinteslas, def:physics:phyx-mini-0626:phyxminiproblems-problemphyxmini0626-chargemagnitudeincoulombs, def:physics:phyx-mini-0626:phyxminiproblems-problemphyxmini0626-numberdensitypercubicmeter, def:physics:phyx-mini-0626:phyxminiproblems-problemphyxmini0626-siliconhallsetup}
  Magnitude form of the Hall-effect law V\_H q n t = I B for a rectangular strip. The width is part of the physical setup but cancels when current density is converted back to total current
\end{definition}
\begin{proof}
  This is the declared model definition of Satisfies Hall Effect Law.
\end{proof}
\begin{definition}[Satisfies Silicon Atomic Density Law]
  \label{def:physics:phyx-mini-0626:phyxminiproblems-problemphyxmini0626-satisfiessiliconatomicdensitylaw}
  \lean{PhyXMiniProblems.ProblemPhyXMini0626.SatisfiesSiliconAtomicDensityLaw}
  \uses{def:physics:phyx-mini-0626:phyxminiproblems-problemphyxmini0626-massinkilograms, def:physics:phyx-mini-0626:phyxminiproblems-problemphyxmini0626-massdensityinkilogramspercubicmeter, def:physics:phyx-mini-0626:phyxminiproblems-problemphyxmini0626-numberdensitypercubicmeter, def:physics:phyx-mini-0626:phyxminiproblems-problemphyxmini0626-siliconhallsetup}
  The mass of one silicon atom is molar mass divided by Avogadro's constant, and bulk mass density is atom number density times mass per atom
\end{definition}
\begin{proof}
  This is the declared model definition of Satisfies Silicon Atomic Density Law.
\end{proof}
\begin{definition}[Satisfies Carrier Per Atom Balance]
  \label{def:physics:phyx-mini-0626:phyxminiproblems-problemphyxmini0626-satisfiescarrierperatombalance}
  \lean{PhyXMiniProblems.ProblemPhyXMini0626.SatisfiesCarrierPerAtomBalance}
  \uses{def:physics:phyx-mini-0626:phyxminiproblems-problemphyxmini0626-numberdensitypercubicmeter, def:physics:phyx-mini-0626:phyxminiproblems-problemphyxmini0626-siliconhallsetup}
  The dimensionless number of conduction electrons per atom converts silicon- atom number density to conduction-electron number density. This is a general balance law and does not assert the requested numerical answer
\end{definition}
\begin{proof}
  This is the declared model definition of Satisfies Carrier Per Atom Balance.
\end{proof}
\begin{definition}[Answer Choice]
  \label{def:physics:phyx-mini-0626:phyxminiproblems-problemphyxmini0626-answerchoice}
  \lean{PhyXMiniProblems.ProblemPhyXMini0626.AnswerChoice}
  The four displayed answer choices
\end{definition}
\begin{proof}
  This is the declared model definition of Answer Choice.
\end{proof}
\begin{definition}[displayed Value]
  \label{def:physics:phyx-mini-0626:phyxminiproblems-problemphyxmini0626-answerchoice-displayedvalue}
  \lean{PhyXMiniProblems.ProblemPhyXMini0626.AnswerChoice.displayedValue}
  \uses{def:physics:phyx-mini-0626:phyxminiproblems-problemphyxmini0626-answerchoice}
  Dimensionless displayed value of an answer choice
\end{definition}
\begin{proof}
  This is the declared model definition of displayed Value.
\end{proof}
\begin{definition}[Is Closest Answer]
  \label{def:physics:phyx-mini-0626:phyxminiproblems-problemphyxmini0626-isclosestanswer}
  \lean{PhyXMiniProblems.ProblemPhyXMini0626.IsClosestAnswer}
  \uses{def:physics:phyx-mini-0626:phyxminiproblems-problemphyxmini0626-answerchoice}
  A choice is at least as close to the physical value as every other choice
\end{definition}
\begin{proof}
  This is the declared model definition of Is Closest Answer.
\end{proof}
\begin{definition}[Rounds To Displayed Value]
  \label{def:physics:phyx-mini-0626:phyxminiproblems-problemphyxmini0626-roundstodisplayedvalue}
  \lean{PhyXMiniProblems.ProblemPhyXMini0626.RoundsToDisplayedValue}
  A value rounds to a displayed center at the given decimal step
\end{definition}
\begin{proof}
  This is the declared model definition of Rounds To Displayed Value.
\end{proof}
% --- Archon declaration topology end ---
% --- Archon physics formalization source end ---
